\chapter{Összegzés}

\section{A projekt összefoglalása}

A dolgozatban bemutatott rendszer mára egy egyszerű premade gráfötletből olyan összetett, de világosan tagolt projektté fejlődött, amely képes:

\begin{itemize}
    \item valódi League of Legends meccsadatokat begyűjteni és adatkészletek szerint kezelni;
    \item a játékosokat \texttt{puuid} alapon normalizálni és szerepkörérzékeny mutatókkal leírni;
    \item az ismétlődő együtt- és egymás ellen játszásokból súlyozott asszociatív gráfokat építeni;
    \item közösségszerkezetet tárolni és újrafelhasználni;
    \item a Rust futtatókörnyezetben útvonalakat, globális nézeteket és analitikákat számolni;
    \item ezeket a kimeneteket több különálló, de koherens frontend oldalon bemutatni.
\end{itemize}

A jelenlegi rendszerállapot alapján a dolgozat legfontosabb üzenete számomra az, hogy ez a projekt már nem csak ötletszinten érdekes. A \texttt{flexset} core-periphery szerkezete, a Flex/SoloQ összehasonlítás, az asszortativitás és a Rust-alapú központisági réteg azt mutatják, hogy a hálózatban ténylegesen megfogható, mérhető szerkezetek jelennek meg. Közben a teljes architektúra is azt támasztja alá, hogy a rendszer nem véletlenszerűen összehordott komponensekből áll, hanem egyre inkább tudatosan szétválasztott felelősségi körök szerint fejlődött.

Az is fontos tanulság, hogy a projekt legerősebb része nem egyetlen „nagy trükk”. Nem az a lényege, hogy van benne 3D nézet, vagy hogy Rustban fut néhány algoritmus, vagy hogy vannak klaszterek. A valódi érték abban van, hogy az adatgyűjtés, a pontszámítás, a gráfépítés, a perzisztencia, az útvonalkeresés és a kutatási analitika végre ugyanannak a rendszernek a része.

\section{A dolgozat védhető fő következtetései}

Az aktuális implementáció és az elvégzett empirikus elemzések alapján a következő állítások tekinthetők védhetőnek:

\begin{itemize}
    \item A League of Legends ismétlődő mérkőzésadataiból értelmes, több rétegben elemezhető játékosháló építhető; ez a flexset dataseten 4980 meccs feldolgozásával valósult meg.
    \item A szerepkörérzékeny, percentilis-alapú \texttt{opscore}/\texttt{feedscore} pipeline kellően értelmezhető ahhoz, hogy gráfanalitikai bemenetként használjuk; a teljes metrikafedettség (5741/5741 a flexseten, 1538/1538 a soloq-on) ezt megerősíti.
    \item A Python és Rust közötti munkamegosztás a jelenlegi projektméretben indokolt és szakmailag jól védhető; a kiegészítő C\# / ASP.NET Core böngésző az adatbázis-engineering alprojekt részeként önállóan is védhető eszközként jelenik meg.
    \item A \texttt{flexset} dataseten az \texttt{opscore} mindkét gráfmódban pozitív asszortativitást mutat (\texttt{social-path}: 0.2346, \texttt{battle-path}: 0.1461); ez a \texttt{soloq} dataseten is teljesül, de gyengébb jel mellett (\texttt{social-path}: 0.1750, \texttt{battle-path}: 0.1625).
    \item A \texttt{flexset} Graph V2 artifact core-periphery asszociatív szerkezetet mutat: hídgazdag központi ally-csoportok és sok kicsi, gyengén kapcsolt periférikus csoport jelenik meg.
    \item A kétdataset-összehasonlítás legszilárdabb eredménye: a Flex Queue lényegesen sűrűbb és ismétlődőbb asszociatív gráfot generál, mint a Solo/Duo Queue; ez az átlagos élsúlyban, klaszterméretben és az asszortativitási különbségekben is megjelenik.
    \item A \texttt{feedscore} \texttt{battle-path} módban közel nulla vagy negatív asszortativitást mutat mindkét dataseten, ami megerősíti, hogy a gráfprojekció megválasztása módszertanilag lényeges, nem puszta technikai beállítás.
\end{itemize}
